% =============================================================================
\section{Related Work}\label{sec:related_work}
% =============================================================================

The closest precedent is Bertsimas et al.~\cite{bertsimas2022where}, who
couple an extended DELPHI compartmental model to a bilinear MIP for
vaccination-site location across US states, solving by a coordinate
descent that alternates discretised dynamics with an LP approximation
of bilinear constraints. Their pipeline is the methodological parallel
to ours; we differ in domain (critical-care surge), coupling
(decoupled forecast-then-optimise rather than iterative joint
optimisation), and objective (decision-aware multi-quantile pinball
rather than MSE on DELPHI states). Our framework also draws on
Elmachtoub and Grigas's Smart Predict-then-Optimize
framework~\cite{elmachtoub2022smart} and on Luo and
Stellato's~\cite{luo2024equitable} neural-ODE-based opioid-facility
location, which embeds discretised dynamics into a McCormick-relaxed
MIP. Multi-quantile training follows the COVID-19 Forecast Hub's
Weighted Interval Score evaluation.

On the forecasting side, statistical and recurrent baselines, including
ARIMA and LSTM/GRU models, remain
competitive when applied per region~\cite{ajao2024deep}. Valente et
al.~\cite{valente2022forecasting} forecast ICU occupancy in Lazio with
a bidirectional LSTM augmented by lagged mobility features and
MC-dropout prediction intervals (MAE $\approx 7$ beds, weekly horizon).
STAN~\cite{gao2021stan} pairs graph attention with an SIR-consistency
loss for case prediction; HOIST~\cite{gao2023hoist} adds an
Ising-dynamics-inspired regulariser over an adaptively learned location
graph for hospitalisation forecasting across 2{,}299 US counties. PINNs
embed compartmental ODEs into the training
loss~\cite{ajao2025hybrid,qian2025physics}. NHS-specific work is
sparse, with the authors' prior
work~\cite{ajao2024deep,ajao2025hybrid} forecasting national MV demand
and existing trust-level ICU studies remaining single-site. None
targets multi-region NHS critical-care planning, and all optimise
statistical rather than operationally aligned metrics.

On the optimisation side, Eriskin et al.~\cite{eriskin2022robust}
present a robust multi-objective MILP for emergency-facility location;
Wan et al.~\cite{wan2023two} solve a two-stage NSGA-II ICU allocation
on Lyon data; Shams Eddin and El Hajj~\cite{shamseddin2025optimization}
integrate SEIRD-driven demand with standard, buffer, and field-hospital
beds plus inter-hospital transfers. In vaccine supply chains,
Kargar et al.~\cite{kargar2024data} couple an agent-based simulator to
a bi-objective MILP, and Tang et al.~\cite{tang2025mobile} apply ALNS
to mobile-vaccination scheduling. Metaheuristics such as
NSGA-II~\cite{deb2002fast} are routinely used at scales where exact
MILP becomes infeasible. The predict-then-optimise framing of
Elmachtoub and Grigas~\cite{elmachtoub2022smart} and the neural-ODE
analogue of Luo and Stellato~\cite{luo2024equitable} for
McCormick-relaxed facility-location optimisation are direct
methodological cousins. None of this literature targets comparative
MILP-and-metaheuristic allocation across the seven NHS England regions
under a decision-aware forecasting objective.
